% !PRINTED 145
\weberpage{145}
%<<CONT
so sind die Coëfficienten $S_0,S_1\ldots S_\mu$ ganze \gesperrt{symmetrische} Functionen der $n-1$ Veränderlichen $\alpha_2,\ldots\alpha_n$.

Bezeichnen wir die symmetrischen Grundfunctionen dieser letzteren Variablen mit $a'_1,a'_2\ldots a'_{n-1}$, so können wir die Coëfficienten $S_0,S_1\ldots S_\mu$ nach unserer Voraussetzung rational durch diese ausdrücken. Es ist aber nach \S.\,42, (2) und (3)
\[
\begin{array}{l}
 a'_1=f_1(\alpha_1)=\alpha_1+a_1\\
 a'_2=f_2(\alpha_1)=\alpha_1^2+a_1\alpha_1+a_2\\
 a'_3=f_3(\alpha_1)=\alpha_1^3+a_1\alpha_1^2+a_2\alpha_1+a_3\\
 \makebox[23em][l]{\dotfill},
\end{array}
\]
d.\ h. die Coëfficienten $S_0,S_1\ldots S_\mu$ können ganz und rational durch $\alpha_1,a_1,a_2\ldots a_n$ ausgedrückt werden. Wenn wir also, nachdem diese Ausdrücke eingeführt sind, in (2) aufs Neue nach Potenzen von $\alpha_1$ ordnen, so erhalten wir
\weq{3}{S=A_0\alpha_1^m+A_1\alpha_1^{m-1}+\cdots+A_{m-1}\alpha_1+A_m,}
worin im Allgemeinen $m$ ein von $\mu$ verschiedener Exponent sein wird, und die Coëfficienten $A_0,A_1\ldots A_m$ ganz und rational von $a_1,a_2\ldots a_n$ abhängen. Wir setzen wieder
\weq{4}{\Phi(x)=A_0x^m+A_1x^{m-1}+\cdots+A_{m-1}x+A_m}
und dividiren $\Phi(x)$ durch
\[
\begin{aligned}
 f(x)&=x^n+a_1x^{n-1}+a_2x^{n-2}+\cdots+a_n\\
     &=(x-\alpha_1)(x-\alpha_2)\ldots(x-\alpha_n).
\end{aligned}
\]
Es ergiebt sich ein Quotient und ein Rest, der in Bezug auf $x$ höchstens vom Grade $n-1$ ist. Wir setzen also
\weq{5}{\Phi(x)=Qf(x)+\psi(x),}
\weq{6}{\psi(x)=C_0x^{n-1}+C_1x^{n-2}+\cdots+C_{n-2}x+C_{n-1},}
und hierin sind $C_0,C_1\ldots C_{n-1}$ ganze rationale Functionen von $a_1,a_2\ldots a_n$, in denen, wenn $S$ in seiner ursprünglichen Form keine gebrochenen Coëfficienten enthält, auch keine Brüche vorkommen.

Nun ist aber, da $f(\alpha_1)$ verschwindet, und $S=\Phi(\alpha_1)$ ist, nach (5)
\weq{7}{S=\psi(\alpha_1)}
und hierin kann, da $S$ symmetrisch ist, $\alpha_1$ durch $\alpha_2,\alpha_3\ldots\alpha_n$ ersetzt werden. Hieraus ergeben sich die folgenden $n$ Gleichungen:
%CONT>>
